\documentclass[11pt,a4paper]{article}
\usepackage[margin=2.5cm]{geometry}
\usepackage{amsmath,amsthm,amssymb}
\usepackage{graphicx}
\usepackage{booktabs}
\usepackage[hidelinks]{hyperref}
\usepackage{cleveref}
\usepackage[backend=biber,style=numeric,urldate=long]{biblatex}
\addbibresource{references.bib}
\setcounter{biburlnumpenalty}{100}
\setcounter{biburlucpenalty}{100}
\setcounter{biburllcpenalty}{100}
\graphicspath{{figures/}}

\theoremstyle{definition}
\newtheorem{definition}{Definition}[section]
\newtheorem{example}[definition]{Example}
\theoremstyle{plain}
\newtheorem{theorem}[definition]{Theorem}
\newtheorem{lemma}[definition]{Lemma}
\theoremstyle{remark}
\newtheorem{remark}[definition]{Remark}

\title{Walks, paths, and connectivity from the adjacency matrix}
\author{}
\date{November 2026}

\begin{document}
\maketitle

\begin{abstract}
For a finite simple undirected graph with vertices $v_1,\dots,v_N$ and adjacency matrix $A$, the entry $(A^k)_{ij}$ is the number of walks of length $k$ from $v_i$ to $v_j$ (walks, not paths; a walk is a sequence of vertices in which each is adjacent to the next). A path is a walk with no repeated vertex. For $N\ge2$, the graph is connected if and only if every off-diagonal entry of $A+A^2+\dots+A^{N-1}$ is positive. Each statement is proved in full and the test is run on three small graphs.
\end{abstract}

\section{Introduction}\label{sec:intro}

A graph is a finite set of points, called vertices, together with lines, called edges, each joining two vertices. The adjacency matrix of a graph is a table of zeros and ones that records which pairs of vertices are joined. This note takes a graph with $N$ vertices and adjacency matrix $A$ (defined in \cref{sec:defs}) and answers three questions, using only matrix multiplication and induction.
\begin{enumerate}
\item[(Q1)] What does the entry $(A^k)_{ij}$ of the $k$-th power of $A$ count? \emph{Answer:} the number of walks of length $k$ from vertex $v_i$ to vertex $v_j$; walks may repeat vertices, so this is not a count of paths (\cref{thm:counting}).
\item[(Q2)] What is a path between two vertices? \emph{Answer:} a walk between them in which no vertex occurs twice (\cref{def:path}).
\item[(Q3)] How can one decide from $A$ alone whether the graph is connected? \emph{Answer:} for $N\ge 2$, the graph is connected if and only if every entry off the main diagonal of $A+A^2+\dots+A^{N-1}$ is positive (\cref{thm:sum}).
\end{enumerate}
The questions matter because the answer to Q3 turns a question about a picture into a computation on a table of numbers. The answers are known; the walk-counting statement appears in several sets of notes \cite{galwaynotes,gotti18211,bai2025adjacency}, the lemma that a walk contains a path in a graph theory text \cite{grinberg2025graph}, and the sum test in a journal article \cite{paramadevan2021adjacency}. Their proofs are written out in full here, because the sources sketch or omit some steps. The support for the answers is these proofs and the examples in \cref{sec:example}, which can be checked by hand; the limits of the test are stated at the end of that section. \Cref{sec:defs} gives the definitions, including the answer to Q2; \cref{sec:counting} answers Q1; \cref{sec:test} answers Q3, which needs the other two.

\section{Definitions}\label{sec:defs}

Throughout, $G$ is a finite, simple, undirected graph: its vertices are $v_1,\dots,v_N$ with $N\ge 1$; no edge joins a vertex to itself; two vertices are joined by at most one edge; and the edge joining $u$ and $v$ is the same as the edge joining $v$ and $u$ \cite[Definition~2.1.1, p.~9]{grinberg2025graph}. Two vertices are \emph{adjacent} if an edge joins them. The indices $i,j,r$ run over $1,\dots,N$.

\begin{definition}[Adjacency matrix]\label{def:matrix}
The \emph{adjacency matrix} of $G$ is the $N\times N$ matrix $A$ with $A_{ij}=1$ if $v_i$ and $v_j$ are adjacent and $A_{ij}=0$ otherwise \cite[Definition~1.1.2]{galwaynotes}. The entries $A_{ii}$ form the \emph{diagonal}; the others are \emph{off-diagonal}.
\end{definition}
Since $G$ is simple and undirected, $A$ is symmetric and its diagonal is zero \cite[p.~3]{galwaynotes}. Powers are the usual ones: $A^0=I$ is the identity matrix, $A^{k+1}=A^kA$, and the entry in row $i$ and column $j$ of a product is $(XY)_{ij}=\sum_{r=1}^{N}X_{ir}Y_{rj}$. All entries of $A$ are non-negative integers, so all entries of every $A^k$ are too.

\begin{definition}[Walk and length]\label{def:walk}
A \emph{walk} from $u$ to $w$ is a sequence of vertices $u_0u_1\cdots u_k$ with $k\ge 0$, $u_0=u$, $u_k=w$, in which $u_{t-1}$ and $u_t$ are adjacent for $1\le t\le k$. Its \emph{length} is $k$. Vertices may repeat, and a walk may traverse the same edge several times, each time counted.
\end{definition}
A single vertex $u$ is a walk of length $0$ from $u$ to itself. This definition, including $k\ge0$, is Definition~2.9.1(a) and (b) of Grinberg \cite[p.~36]{grinberg2025graph}.

\begin{definition}[Path, path between two vertices]\label{def:path}
A \emph{path} is a walk whose vertices are pairwise distinct \cite[Definition~2.9.1(c), p.~36]{grinberg2025graph}. A \emph{path between $u$ and $w$} (or \emph{from $u$ to $w$}) is a path $u_0\cdots u_k$ with $u_0=u$ and $u_k=w$. If $u\ne w$ it has length at least $1$; the only path from $u$ to itself is the single vertex $u$, because $u_0=u_k$ with $k\ge1$ would repeat a vertex.
\end{definition}

\begin{definition}[Connected]\label{def:connected}
$G$ is \emph{connected} if for any two vertices $u,w$ there is a path from $u$ to $w$.
\end{definition}
Grinberg defines connectedness through components; by his Proposition~2.9.10 \cite[p.~42]{grinberg2025graph} the two definitions agree. A vertex is joined to itself by the path of length $0$, so only pairs of different vertices matter.

\section{Q1: what the entries of \texorpdfstring{$A^k$}{Ak} count}\label{sec:counting}

\begin{theorem}[Walk counting]\label{thm:counting}
For every integer $k\ge 0$ and all $i,j$, the entry $(A^k)_{ij}$ is the number of walks of length $k$ from $v_i$ to $v_j$.
\end{theorem}
The sources prove this starting from $k=1$ \cite[Theorem~1.1.3, p.~4]{galwaynotes}\cite[Proposition~2, p.~1]{gotti18211}\cite[Theorem~1, p.~1]{bai2025adjacency}; the case $k=0$ is covered below.
\begin{proof}
By induction on $k$.

\emph{Base case $k=0$.} $A^0=I$. A walk of length $0$ from $v_i$ to $v_j$ is a single vertex that is both start and end, so there is exactly one such walk if $i=j$ and none if $i\ne j$. This is $I_{ij}$.

\emph{Inductive step.} Assume that for some $k\ge0$ and all $i,r$ the entry $(A^k)_{ir}$ is the number of walks of length $k$ from $v_i$ to $v_r$. Fix $i$ and $j$. Since $A^{k+1}=A^kA$,
\begin{equation}\label{eq:step}
(A^{k+1})_{ij}=\sum_{r=1}^{N}(A^k)_{ir}A_{rj}.
\end{equation}
Write a walk of length $k+1$ from $v_i$ to $v_j$ as $w_0w_1\cdots w_kw_{k+1}$ with $w_0=v_i$ and $w_{k+1}=v_j$; its last vertex is $w_{k+1}$ and the vertex before the last is $w_k$. Exactly one $r$ has $w_k=v_r$, so these walks fall into $N$ classes $\mathcal C_r$, one for each $r$: $\mathcal C_r$ is the set of walks of length $k+1$ from $v_i$ to $v_j$ whose vertex before the last is $v_r$. Every such walk lies in exactly one class.

Fix $r$. If $A_{rj}=1$, that is, $v_r$ and $v_j$ are adjacent, then deleting the final vertex $v_j$ from a walk in $\mathcal C_r$ leaves a walk of length $k$ from $v_i$ to $v_r$; and appending $v_j$ to a walk of length $k$ from $v_i$ to $v_r$ gives a walk in $\mathcal C_r$, because $v_r$ is adjacent to $v_j$. The two operations undo each other, so $\mathcal C_r$ and the set of walks of length $k$ from $v_i$ to $v_r$ are in one to one correspondence. By the hypothesis, the number $|\mathcal C_r|$ of walks in $\mathcal C_r$ is $(A^k)_{ir}=(A^k)_{ir}A_{rj}$. If $A_{rj}=0$, then $v_r$ and $v_j$ are not adjacent, so no walk can have $w_k=v_r$ and $w_{k+1}=v_j$, and $|\mathcal C_r|=0=(A^k)_{ir}A_{rj}$.

Adding over the $N$ classes, the number of walks of length $k+1$ from $v_i$ to $v_j$ is $\sum_r(A^k)_{ir}A_{rj}$, which is $(A^{k+1})_{ij}$ by \cref{eq:step}.
\end{proof}

\paragraph{The entries count walks, not paths.} The proof uses only that consecutive vertices are adjacent, so walks with repeated vertices are counted. Take the graph $G_1$ with $N=3$ and edges $v_1v_2$ and $v_2v_3$. Its matrix and the next two powers are
\[
A=\begin{pmatrix}0&1&0\\1&0&1\\0&1&0\end{pmatrix},\quad
A^2=\begin{pmatrix}1&0&1\\0&2&0\\1&0&1\end{pmatrix},\quad
A^3=\begin{pmatrix}0&2&0\\2&0&2\\0&2&0\end{pmatrix}.
\]
By \cref{thm:counting}, $(A^2)_{11}=1$ counts the one walk $v_1v_2v_1$ of length $2$ from $v_1$ to $v_1$. It repeats $v_1$, so it is not a path; the only path from $v_1$ to $v_1$ is $v_1$ alone, of length $0$ (\cref{def:path}). So ``$(A^2)_{11}$ is the number of paths of length $2$ from $v_1$ to $v_1$'' is false, although ``the number of walks'' is true. Reading the entry as ``$1$ if $v_j$ can be reached from $v_i$'' is also wrong: $(A^3)_{12}=2$ counts the walks $v_1v_2v_1v_2$ and $v_1v_2v_3v_2$, and $(A^3)_{13}=0$ although a path $v_1v_2v_3$ joins $v_1$ to $v_3$ (it has length $2$, not $3$).

\section{Q3: a connectivity test from \texorpdfstring{$A$}{A}}\label{sec:test}

The test rests on two lemmas about walks and paths. \Cref{lem:cut} is Corollary~2.9.7 of Grinberg \cite[p.~39]{grinberg2025graph}; \cref{lem:reach} combines the walk-counting theorem with the bound that a path has at most $N-1$ edges (it has at most $N$ distinct vertices), which Grinberg states in Section~2.9.4 \cite[p.~40]{grinberg2025graph}. The proofs below are written out in full.

\begin{lemma}[Every walk contains a path]\label{lem:cut}
Let $W=u_0\cdots u_k$ be a walk from $u$ to $w$. Then there is a path from $u$ to $w$ of length at most $k$.
\end{lemma}
\begin{proof}
By strong induction on $k$ (strong, because $W'$ below can have any length from $0$ to $k-1$): for a walk of length $k$ we assume the statement for all walks of smaller length.

\emph{Base case $k=0$.} $W=(u)$ has one vertex, so it is a path of length $0$.

\emph{Inductive step, $k\ge1$.} If the vertices of $W$ are pairwise distinct, $W$ is itself a path of length $k$. Otherwise there are positions $s<t$ with $u_s=u_t$. Delete the vertices at positions $s+1,\dots,t$; this removes the second copy $u_t$ of the repeated vertex and everything between the two copies, and gives $W'=u_0\cdots u_su_{t+1}\cdots u_k$, and $W'=u_0\cdots u_s$ if $t=k$. $W'$ is a walk, because its only new consecutive pair is $u_s,u_{t+1}$, and $u_s=u_t$ is adjacent to $u_{t+1}$. It starts at $u_0=u$, ends at $u_k=w$ (if $t=k$ it ends at $u_s=u_t=u_k$), and has length $k-(t-s)<k$. The induction hypothesis, applied to the shorter walk $W'$, gives a path from $u$ to $w$ of length at most $k-(t-s)\le k$.
\end{proof}

\begin{lemma}[Reachability from the powers]\label{lem:reach}
Let $i\ne j$. There is a path from $v_i$ to $v_j$ if and only if $(A^k)_{ij}>0$ for some $k$ with $1\le k\le N-1$.
\end{lemma}
\begin{proof}
($\Rightarrow$) Let a path from $v_i$ to $v_j$ have length $k$. As $i\ne j$, $k\ge1$ (a path of length $0$ has a single vertex, so it cannot join different vertices). A path of length $k$ has $k+1$ pairwise distinct vertices among the $N$ vertices of the graph, so $k+1\le N$, that is $k\le N-1$. A path is a walk, so by \cref{thm:counting} $(A^k)_{ij}\ge1$.

($\Leftarrow$) If $(A^k)_{ij}>0$ with $1\le k\le N-1$, then by \cref{thm:counting} a walk of length $k$ joins $v_i$ to $v_j$, and by \cref{lem:cut} a path joins them.
\end{proof}

\begin{theorem}[Sum test]\label{thm:sum}
Let $N\ge2$ and $B=A+A^2+\dots+A^{N-1}$. Then $G$ is connected if and only if $B_{ij}>0$ for all $i\ne j$.
\end{theorem}
For $N=2$ the sum has the single term $B=A$.
For $N=1$ there is no pair of different vertices and the graph is connected. \Cref{thm:sum} is Theorem~2 of Paramadevan and Sotheeswaran \cite[p.~17]{paramadevan2021adjacency}.
\begin{proof}
Every $A^k$ has non-negative entries, so $B_{ij}=\sum_{k=1}^{N-1}(A^k)_{ij}$ is a sum of non-negative terms, and it is positive exactly when at least one term is positive. For $i\ne j$ this holds exactly when a path joins $v_i$ and $v_j$, by \cref{lem:reach}. So $B_{ij}>0$ for all $i\ne j$ exactly when every pair of different vertices is joined by a path, which is \cref{def:connected}.
\end{proof}
The cutoff $N-1$ is the largest length a path can have (proof of \cref{lem:reach}). The entries of $B$ are walk counts, not path counts; only whether they are zero is used. The test can also be written as ``$(I+A)^{N-1}$ has no zero entry'' \cite[Proposition~3, p.~2]{gotti18211}\cite[Theorem~2, p.~1]{bai2025adjacency}. Indeed $(I+A)^{N-1}=\sum_{k=0}^{N-1}\binom{N-1}{k}A^k$ because $I$ and $A$ commute; every coefficient is at least $1$, the term $k=0$ makes the diagonal positive, and off the diagonal only the terms $k\ge1$ contribute, each a positive multiple of $(A^k)_{ij}$ (every coefficient is at least $1$), so the entry is positive exactly when some $(A^k)_{ij}$ with $1\le k\le N-1$ is positive, that is, exactly when $B_{ij}>0$.

\section{Examples and limits}\label{sec:example}

\paragraph{A connected graph.} Take $G_1$ of \cref{sec:counting} ($N=3$, edges $v_1v_2$, $v_2v_3$), whose $A$ and $A^2$ are printed there. Then
\[
B=A+A^2=\begin{pmatrix}1&1&1\\1&2&1\\1&1&1\end{pmatrix}.
\]
Every off-diagonal entry of $B$ is positive, so $G_1$ is connected by \cref{thm:sum}. The pairs $v_1v_2$ and $v_2v_3$ get their positive entry from $A$, which counts the one-step walks $v_1v_2$ and $v_2v_3$. The pair $v_1v_3$ gets $B_{13}=0+1=1$ from $(A^2)_{13}=1$, which counts the walk $v_1v_2v_3$; this walk is also a path.

\paragraph{A disconnected graph.} Let $G_2$ have $N=3$, the edge $v_1v_2$, and the isolated vertex $v_3$. Then
\[
A=\begin{pmatrix}0&1&0\\1&0&0\\0&0&0\end{pmatrix},\quad
A^2=\begin{pmatrix}1&0&0\\0&1&0\\0&0&0\end{pmatrix},\quad
B=\begin{pmatrix}1&1&0\\1&1&0\\0&0&0\end{pmatrix}.
\]
The entries $B_{13}$, $B_{23}$, $B_{31}$, $B_{32}$ are $0$, so $G_2$ is not connected by \cref{thm:sum}. This agrees with the graph: $v_3$ has no neighbour, so no walk leaves it.

\paragraph{The sum must reach $A^{N-1}$.} Let $G_3$ have $N=4$ and edges $v_1v_2$, $v_2v_3$, $v_3v_4$. Then
\[
A=\begin{pmatrix}0&1&0&0\\1&0&1&0\\0&1&0&1\\0&0&1&0\end{pmatrix},\;
A^2=\begin{pmatrix}1&0&1&0\\0&2&0&1\\1&0&2&0\\0&1&0&1\end{pmatrix},\;
A^3=\begin{pmatrix}0&2&0&1\\2&0&3&0\\0&3&0&2\\1&0&2&0\end{pmatrix}.
\]
The entry $(1,4)$ is $0$ in $A$ and in $A^2$ (the only neighbour of $v_1$ is $v_2$, which is not adjacent to $v_4$, so no walk of length $2$ ends at $v_4$) and is $1$ in $A^3$, counting the walk $v_1v_2v_3v_4$. A sum stopped at $A^2$ would have $(1,4)$ entry $0$ and so wrongly report $G_3$ as disconnected; the path $v_1v_2v_3v_4$ has length $3=N-1$ (\cref{lem:reach}). The full sum is
\[
B=A+A^2+A^3=\begin{pmatrix}1&3&1&1\\3&2&4&1\\1&4&2&3\\1&1&3&1\end{pmatrix},
\]
whose off-diagonal entries are all positive (the diagonal entries are $B_{ii}$), so $G_3$ is connected by \cref{thm:sum}.

\paragraph{Limits.} The test takes only $A$ as input, and it has the following limits. It reads only whether entries are zero: a positive entry of $A^k$ means a walk of length exactly $k$ exists. Hence (\cref{lem:cut}) a path of length at most $k$ exists, not necessarily one of length exactly $k$ (in $G_1$, whose $A^3$ is printed in \cref{sec:counting}, $(A^3)_{12}=2$, yet the only path from $v_1$ to $v_2$ has length $1$). The diagonal of $B$ is left out on purpose: a connected graph can have zeros there, for example $N=2$ with one edge, where $B=A$ has a zero diagonal, so requiring a positive diagonal would reject a connected graph. The test is proved here only for finite simple undirected graphs with $N\ge2$.


\printbibliography

\end{document}
